
\section{Resumo}
\label{sec:ik_resumo}

A cinemática inversa permite determinar os ângulos articulares do myCobot~280 a partir da posição pretendida $(x,y,z)$ e da inclinação da garra $\psi$.

A resolução segue três etapas: cálculo da orientação da base, determinação da posição do pulso e cálculo dos ângulos das juntas J2, J3 e J4.

Considerando $D=L_6+L_7$, $L_A=L_2$ e $L_B=L_3$, as expressões utilizadas são apresentadas de seguida.

\subsubsection*{1. Plano horizontal}

A partir das coordenadas do alvo, calculam-se a distância ao eixo da base, a distância frontal e a rotação da junta J1:

\begin{align*}
    r &= \sqrt{x^2+y^2}\\
    d &= \sqrt{r^2-L_5^2}\\
    q_1 &= \operatorname{atan2}(y,x)
    -\operatorname{atan2}(d,L_5)
\end{align*}

\subsubsection*{2. Posição do pulso}

Definindo a inclinação da garra, determina-se a posição do pulso relativamente ao ombro, descontando as contribuições dos segmentos $L_4$ e $D$:

\begin{align*}
    \alpha_3 &= -\psi\\
    x_{3d} &= d-D\cos\alpha_3
    -L_4\cos(\alpha_3+90^\circ)\\
    y_{3d} &= (z-L_1)-D\sin\alpha_3
    -L_4\sin(\alpha_3+90^\circ)
\end{align*}

\subsubsection*{3. Braço planar de 2 DOF}

Conhecida a posição do pulso, calculam-se a distância $r_3$ e os ângulos auxiliares através da lei dos cossenos:

\begin{align*}
    r_3 &= \sqrt{x_{3d}^2+y_{3d}^2}\\
    \alpha_1 &= \operatorname{atan2}(y_{3d},x_{3d})\\
    \beta_1 &= \arccos\left(
    \frac{L_A^2+r_3^2-L_B^2}{2L_A r_3}
    \right)\\
    \beta_2 &= \arccos\left(
    \frac{L_A^2+L_B^2-r_3^2}{2L_A L_B}
    \right)
\end{align*}

Para a configuração com o cotovelo para cima, obtêm-se:

\begin{align*}
    \theta_1 &= \alpha_1+\beta_1\\
    \theta_2 &= -(180^\circ-\beta_2)\\
    \theta_3 &= (\alpha_3+90^\circ)-\theta_1-\theta_2
\end{align*}

Finalmente, os ângulos articulares do manipulador são:

\begin{equation}
    \boxed{
    \begin{aligned}
        q_2 &= \theta_1-90^\circ\\
        q_3 &= \theta_2\\
        q_4 &= \theta_3\\
        q_5 &= 0^\circ
    \end{aligned}
    }
    \label{eq:ik_final}
\end{equation}

A junta J6 é utilizada para alinhar os dedos da garra com a orientação $\gamma$ de uma peça cúbica, quando esta é conhecida e $\psi>45^\circ$:

\begin{equation}
    q_6=
    \operatorname{mod}(q_1-\gamma+45^\circ,\;90^\circ)
    -45^\circ
\end{equation}

Caso contrário, considera-se $q_6=0^\circ$.

\subsubsection*{Verificação das soluções}

Uma solução apenas é aceite quando são respeitadas as condições de alcance e as restrições implementadas:

\begin{itemize}
    \item A distância horizontal deve satisfazer $r>L_5$;
    \item O pulso deve respeitar $|L_2-L_3|\leq r_3\leq L_2+L_3$;
    \item Os ângulos articulares devem respeitar os limites permitidos, com $|q_3|\leq140^\circ$, e as juntas J3, J4 e J5 devem permanecer acima do nível de referência.
\end{itemize}

\subsubsection*{Procedimento do algoritmo}

O algoritmo completo segue os seguintes passos:

\begin{enumerate}
    \item Converter a posição do alvo do referencial do mundo para o da base, através da Equação~\eqref{eq:mundo_base};

    \item Calcular $r$, $d$ e $q_1$, rejeitando o alvo se $r\leq L_5$;

    \item Tentar a inclinação preferencial $\psi=90^\circ$, calculando os restantes ângulos e verificando as condições de validade;

    \item Se não existir uma solução válida, testar inclinações alternativas em incrementos de $5^\circ$, no intervalo de $30^\circ$ a $135^\circ$;

    \item Quando é encontrada uma inclinação válida, realizar um refinamento por bisseção para procurar uma orientação mais próxima da preferencial;

    \item Caso nenhuma inclinação permita obter uma configuração válida, considerar o alvo inalcançável segundo as restrições implementadas.
\end{enumerate}

Para uma inclinação fixa, podem existir duas soluções geométricas, correspondentes ao cotovelo para cima e para baixo. Na implementação, é utilizada por defeito a primeira configuração.

A consistência dos resultados é verificada aplicando a cinemática direta, apresentada na Equação~\eqref{eq:dk_final}, aos ângulos obtidos, como demonstrado na Secção~\ref{sec:ik_resultados}.

\clearpage
